\section{Introduction}
\label{sec:intro}

The rapid expansion of specialized deep learning models has made multi-task adaptation a critical yet resource-intensive endeavor. Traditionally, multi-task learning requires joint fine-tuning across all downstream distributions, which is computationally expensive and prone to negative transfer~\cite{jacobs1991adaptive}. To address these limitations, weight-space \textit{model merging} has emerged as an efficient, post-hoc paradigm~\cite{wortsman2022model}. By interpolating the parameters of specialized task-specific networks (i.e., task vectors) back into a single base model, practitioners can combine capabilities without additional training data or parameter overhead~\cite{ilharco2022editing, yadav2023ties}.

While early work utilized static uniform merging, recent research focuses on dynamically adapting merging weights at test-time (Test-Time Adaptation, TTA) using unsupervised objective functions on unlabeled inference streams~\cite{adamerging, polymerge, zipmerge}. Recently, a highly stylized and complex class of ``quantum-inspired'' techniques, epitomized by Quantum Wavefunction Superposition Merging (QWS-Merge)~\cite{qwsmerge}, has claimed state-of-the-art multi-task accuracy. QWS-Merge abandons standard linear routing to model task-specific weights as eigenstates in parameter Hilbert space, representing merging as a coherent quantum-like superposition that collapses based on wave-like phase interference between input features and a trainable layer-wise basis.

As \textbf{The Methodologist}, we approach these ``quantum-inspired'' claims with extreme skepticism. Historically, deep learning has been susceptible to the adoption of elaborate mathematical metaphors—such as physical thermodynamics or quantum mechanics—which often mask simpler underlying statistical mechanisms. A rigorous inspection of QWS-Merge's evaluation protocol reveals two critical flaws:
\begin{enumerate}
    \item \textbf{Crippled Baselines:} The authors compare QWS-Merge against a ``Linear Router'' classical baseline that is intentionally crippled. Specifically, the baseline is global (bypassing the crucial layer-wise parameter specialization) and utilizes an unconstrained high-dimensional parameter projection matrix ($768$ parameters), making it highly susceptible to overfitting.
    \item \textbf{Omission of Basic Regularization:} The catastrophic collapse of the classical Linear Router on out-of-distribution (OOD) tasks (e.g., SVHN) is presented as a fundamental limitation of classical linear representations. However, this classical router was trained in a low-data regime (a tiny $64$-sample calibration split) with \textit{zero regularization} (no weight decay or input normalization). 
\end{enumerate}

In this work, we isolate the true drivers of multi-task routing performance by stripping away the quantum mechanical vocabulary. We propose a simple, transparent classical alternative: the \textbf{Layer-wise Low-dimensional Classical Router (L3-Router)}. Operating in the exact same low-dimensional unit-projected space as QWS-Merge, our L3-Router projects representations linearly to layer-wise task-merging coefficients. We formulate three variants of our classical routing architecture: L3-Linear, L3-Tanh (bounded), and L3-Softmax (normalized). We then train these routers using standard, classical $L_2$ weight decay (regularization) on the calibration set. We explicitly note that our L3-Router achieves a relative \textbf{16.7\%} parameter footprint reduction over QWS-Merge (saving exactly 56 parameters, from 336 down to 280); while this absolute saving is practically negligible relative to backbone model scales (meaning it offers no real-world hardware storage savings), it holds strong theoretical value by proving that classical linear channels can achieve superior dynamic capacity without requiring any auxiliary wave amplitude or phase variables.

We evaluate these formulations within a controlled representation sandbox. While evaluating on a sandbox surrogate introduces a scale gap compared to full-scale Vision Transformers (e.g., ViT-B/16 CLIP) or large language models (LLMs), this controlled environment serves as an essential isolating coordinate space to decouple routing dynamics from weight space misalignment. We acknowledge this scale gap explicitly and provide an immediately actionable deployment roadmap in Appendix~\ref{sec:roadmap_app} to guide practitioners in scale-verifying our findings on CLIP and LLM merges.

Our empirical results within our controlled multi-task representation sandbox reveal a clear picture:
\begin{itemize}
    \item \textbf{Limitations of Quantum-Inspired Formulations:} The wave-interference cosine activation of QWS-Merge is functionally unstable, completely collapsing to a Joint Mean of \textbf{36.10\%} (even worse than uniform merging's \textbf{43.40\%}) and yielding near-random performance of \textbf{2.00\%} on the out-of-distribution SVHN task.
    \item \textbf{Simpler Layer-wise Routing is Highly Stable:} By replacing the cosine wave formulation with our proposed classical \textbf{L3-Linear} router, we avoid this catastrophic collapse and achieve a Joint Mean accuracy of \textbf{63.10\%} (\textbf{+27.00\%} absolute improvement over the quantum SOTA).
    \item \textbf{The Ultimate Baseline Confounder:} Most remarkably, we find that the simplest, global unregularized classical \textbf{Linear Router} baseline outperforms all other models, achieving a Joint Mean accuracy of \textbf{67.20\%}. This proves that complex, layer-wise specialized routing represents an over-engineered mechanism compared to a simple, global classical mapping, highlighting a major blind spot in current model-merging literature.
    \item \textbf{Mitigating Heterogeneity Collapse:} Under realistic deployment streams, mixed-task batches cause a severe \textit{heterogeneity collapse}, dropping Linear Router accuracy from \textbf{67.20\%} to \textbf{51.10\%} (\textbf{-16.10\%}) and QWS-Merge to \textbf{10.80\%} (\textbf{-25.30\%}). However, our proposed \textbf{L3-Softmax} router exhibits extreme robustness to batch mixedness, dropping by only \textbf{4.10\%} and matching the linear router at \textbf{50.30\%}.
\end{itemize}

By demonstrating that a simple, global classical linear baseline beats both complex quantum-inspired model merging and our own multi-layer methods, this paper serves as a vital cautionary tale. We advocate for rigorous baseline tuning, ablation of structural complexity, and strict scientific hygiene when evaluating model-merging frameworks.
